\section{Barridos de parámetros y análisis de sensibilidad}
\label{sec:barridos}

Uno de los requisitos del proyecto es que los filtros también se sintonicen,
como se sintoniza un PID, y que el software barra sus parámetros de forma
automática. El
módulo \cod{quadsim/sweeps.py} (con \cod{objectives.py} y
\cod{experiments/sweep\_cli.py}) lo hace; el informe completo está en
\cod{docs/BARRIDOS.md}. Esta sección resume el método y los hallazgos que
tocan al modelo.

\subsection{Método}

{\sloppy Cada parámetro se nombra con una ruta de texto (\cod{filter.rp.0.fc},
\cod{estimator.ekf.accel\_noise}, \cod{pid.rp.rate\_kd\_scale},
\cod{sensors.vib\_gyro\_dps}, \ldots) que se aplica a la configuración o a la
simulación viva solo con su API pública. Los diseños de experimento son
grilla completa, aleatorio o hipercubo latino \cite{mckay79}, con
refinamiento local alrededor del mejor punto (paso exploratorio de Hooke y
Jeeves \cite{hooke61}). Hay tres ingredientes que vuelven útil la
comparación:\par}

\begin{itemize}
\item \textbf{Números aleatorios comunes} \cite{lawkelton}: la semilla
depende de la réplica y no del punto, así que todos los puntos ven el mismo
ruido, el mismo sesgo y la misma turbulencia; las diferencias entre puntos
vienen de los parámetros, y dos réplicas ya muestran tendencias limpias.
\item \textbf{Costo normalizado}, $J=\sum_kw_k\,m_k/m_k^{base}$, donde la base
es la configuración por defecto con las mismas semillas ($J<1$ es mejor que
el defecto), junto con el frente de Pareto de dos métricas, que es la
respuesta honesta cuando hay compromiso.
\item \textbf{Condiciones}: parámetros que son condiciones de operación
(vibración, EMI) y no ajustes. Cada punto se corre en todas y el ranking es
robusto (promedio o peor caso), para que el ``mejor'' punto no sea siempre
el de la condición más fácil.
\end{itemize}

Los escenarios de evaluación incluyen escalones de actitud y un
\emph{doblete de tasa} de $\pm90^\circ$/s, porque el retardo de un filtro casi
no se nota en el escalón de ángulo (el error de seguimiento varió entre
2.02 y $2.07^\circ$ con cortes de 20 a 200~Hz) y sí en la respuesta de tasa.
El ruido que llega a los motores se mide como el RMS del acelerador por
encima de 20~Hz en vuelo estacionario (\cod{motor\_hf\_pct}): un motor con
$T_m=20$~ms no sigue mandos más rápidos que unos 8~Hz, y esa energía se va
en calor. Es un sustituto, porque el modelo de motor no tiene pérdidas
eléctricas.

\subsection{Hallazgos que tocan al modelo}

\begin{itemize}
\item \textbf{La vibración está donde dice el modelo.} En todas las corridas
el tono dominante del giroscopio crudo está en 177.3--177.6~Hz, la
frecuencia de paso de pala calculada.
\item \textbf{Pasa-bajos: compromiso asimétrico.} Subir el corte de 80 a
200~Hz mejora el sobreimpulso de tasa 3.5 puntos pero multiplica por 13 el
ruido de motor; bajarlo a 20~Hz casi no quita más ruido y dispara el
sobreimpulso al 51~\%. Con el mismo corte, un $Q=1.2$ atrasa menos que
Butterworth.
\item \textbf{Notch.} Centrado en 178~Hz con $Q=1$ reduce el ruido de motor
63~\% con la mitad de retardo, y permite subir el pasa-bajos a 150~Hz. Un
notch angosto ($Q=8$) descentrado 2~Hz pierde casi todo su beneficio: es la
razón de los notch dinámicos ligados a las rpm.
\item \textbf{El derivativo SIMC era bajo.} Con $k_d=K_c\theta/3$ el doblete
de tasa tiene 20~\% de sobreimpulso; con el triple, 5.3~\%, con solo 7~\% más
de ruido de motor. Kd y corte del filtro deben moverse juntos: bajar el
corte a 40--60~Hz y subir $k_d$ de 3 a 4 veces da 3--6~\% de sobreimpulso
con 40--70~\% menos ruido.
\item \textbf{Estimadores.} El EKF con $\sigma_a=\sigma_m=0.05$ estaba mal
sintonizado: sin EMI, $(\sigma_a,\sigma_m)\approx(1.1,\,0.007)$ baja el error
de roll/pitch y de yaw 3 y 7 veces; el punto robusto frente a EMI es
$(0.5,\,0.3)$; $\beta$ de Madgwick óptimo entre
0.10 y 0.15; Mahony por defecto a 3~\% de su óptimo. Con EMI ninguna
ganancia arregla el rumbo: en Madgwick y Mahony el magnetómetro corrige los
tres ejes y contamina roll y pitch.
\item \textbf{Los barridos encuentran los huecos del costo.} Una primera
versión del costo sin el tiempo de establecimiento del ángulo llevó al
optimizador a bajar la ganancia de tasa y ``ganar'' con un dron lento; y
sintonizar el filtro solo en vuelo estacionario con turbulencia lleva a un
filtro excesivo (20~Hz). El escenario y el costo deben excitar la dinámica
que se quiere proteger.
\end{itemize}

\paragraph{Qué transfiere al dron real.} La forma de los compromisos, no los
números: la amplitud y el espectro de la vibración son supuestos. Antes de
transferir una sintonía de filtros hay que medir el espectro real del
giroscopio en hover, ajustar $A_g$ en el simulador hasta que el pico
coincida, y repetir el barrido.

\paragraph{Qué se aplicó al núcleo.} A partir de estos barridos y de la
verificación de este documento se cambiaron los valores por defecto: $k_d$ de
tasa de $K_c\theta/3$ a $K_c\theta$, pasa-bajos del giroscopio de 80 a 60~Hz,
ruidos del EKF a $(\sigma_a,\sigma_m)=(0.5,0.3)$, Butterworth de 30~Hz sobre el
acelerómetro antes del AHRS y estimador por defecto Mahony. Para $\beta$ de
Madgwick \emph{no} se tomó el óptimo del barrido (0.10--0.15) sino 0.015: el
barrido corría el estimador en sombra con el acelerómetro crudo, y en lazo
cerrado con el acelerómetro filtrado un $\beta$ grande produce un ciclo
límite (\cref{sec:estimadores}). Con los nuevos valores, el escenario
\cod{attitude\_step} pasó de 20.1 a 5.9~\% de sobreimpulso de tasa y de
0.78 a 0.43$^\circ$ de error de estimación con Madgwick (0.51$^\circ$ con
Mahony); en \cod{estimator\_accuracy} el error RMS bajó de 0.70 a 0.63$^\circ$
con Mahony y de 1.14 a 0.68$^\circ$ con el EKF, cuyo error de rumbo pasó de
2.55 a 0.43$^\circ$. Las tablas de \cod{docs/BARRIDOS.md} se produjeron con
el núcleo anterior.

\paragraph{Barridos de estimadores: sombra y vuelo.} Los barridos de
\cod{madgwick\_beta}, \cod{mahony\_gains} y \cod{ekf\_noise} corren el
estimador candidato en sombra. Hasta la revisión de octubre de 2026 esa
cadena no era la que se vuela: recibía el acelerómetro sin el Butterworth
previo al AHRS y una aceleración cinemática calculada por un filtro de
navegación alimentado con el estimador ideal, sin el lazo autorreferente de
la \cref{sec:leishman}. La revisión externa midió 1.26$^\circ$ RMS en sombra
frente a 0.43$^\circ$ en el lazo para Madgwick. Hoy la sombra y el vuelo usan
la misma función (\cod{quadsim/estimation.py}); los resultados de esos
barridos en \cod{docs/BARRIDOS.md} se obtuvieron con la cadena anterior, y
por eso el $\beta$ por defecto se decidió con mediciones en lazo cerrado y
no con el óptimo del barrido.
